\section{Conclusion}
\label{sec:conclusion}

This paper reorganizes the repository's security work into a research pipeline that begins with public incidents and ends with replay-oriented exploit claims. The central empirical result is straightforward: raw incident volume is much broader than confirmed pure on-chain exploitability. Within the current artifact set, the normalized corpus contains 1{,}378 DeFi-labeled incidents, but the 24-protocol verification study identifies only one high-confidence confirmed pure on-chain exploit path and three additional conditional cases. Replay engineering narrows the lens further, with concrete logic currently implemented for three representative historical exploits.

The paper's stronger contribution is methodological. By separating incident evidence, protocol evidence, and replay evidence, the refactored workflow makes it easier to write papers that are simultaneously richer and more disciplined. It encourages stronger claims where the evidence is strong and more restrained language where the repository still contains only a scaffold, a shortlist, or a conditional finding. That is a better fit for DeFi security research than either a flat incident catalog or a scanner-driven vulnerability count.

The immediate next milestone is to convert the replay layer from implemented logic into archive-validated results. Once that happens for a broader set of cases, the same paper system can support a substantially stronger version of this study with deeper appendices, more complete case tables, and sharper comparative claims. In that sense, the present draft should be read as a credible research baseline: its main value is that it makes the next increment of evidence accumulation straightforward rather than ad hoc.
